\chapter{Libreria \texttt{pgm-core}}
%--------1---------2---------3---------4---------5---------6---------7---------8
Questa libreria contiene delle interfaccie generalizzate per gestire le origine
dati e destinazione dati che possono essere utili nella elaborazione di un
programma.

\VerbatimFootnotes
\section{Sorgente dati --- \texttt{SourceResource}} \label{sec:srcRes}
%--------1---------2---------3---------4---------5---------6---------7---------8
Una generica sorgente dati è una classe che implementa l'\,interfaccia
\texttt{SourceResource}, vedi lis.~\ref{lst:srcRef}, dove il metodo \texttt{get}
fornisce la risorsa \texttt{U} che può essere usata in un
\texttt{try-with-resources}, e i metodi \texttt{iterator} e \texttt{supplier}
permettono di \textsl{leggere} gli elementi mediante un loop\footnote{%
    uso di una \texttt{SourceResource} con un \texttt{Iterator}:
    \begin{javacode}
        try (U u = source.get()) {
        Iterator<I> iterator = source.iterator(u);
        while (iterator.hasNext()) {
            I i = iterator.next();
            // TODO
        }
    }
    \end{javacode}
} o dei supplier\footnote{%
    uso di una \texttt{SourceResource} con un \texttt{Supplier}:
    \begin{javacode}
        try (U u = source.get()) {
        Supplier<I> supplier = source.supplier(u);
        I i = supplier.get();
        while (i != null) {
            // TODO
            i = supplier.get();
        }
    }
    \end{javacode}
}.
%--------1---------2---------3---------4---------5---------6---------7---------8

%--------1---------2---------3---------4---------5---------6---------7---------8
\begin{elisting}[!htb]
\begin{javacode}
public interface SourceResource<U extends AutoCloseable, I> {
    U get();
    Iterator<I> iterator(U u);
    Supplier<I> supplier(U u);
}
\end{javacode}
\caption{Metodi per aprire la risorsa e leggere i dati}
\label{lst:srcRef}
\end{elisting}
%--------1---------2---------3---------4---------5---------6---------7---------8

%--------1---------2---------3---------4---------5---------6---------7---------8
L'\,interfaccia oltre a indicare i metodi di istanza offre dei costruttori
statici facilitare l'\,implementazione della interfaccia.
%--------1---------2---------3---------4---------5---------6---------7---------8

%--------1---------2---------3---------4---------5---------6---------7---------8
\begin{elisting}[!htb]
\begin{javacode}
public interface SourceResource<U extends AutoCloseable, I> {
    static SourceResource<BufferedReader, String> bufferedReader(File file, Charset cs, Runnable inc);
    static SourceResource<BufferedReader, String> bufferedReader(File file, Runnable inc);
    static SourceResource<BufferedReader, String> bufferedReader(File file, Charset cs);
    static SourceResource<BufferedReader, String> bufferedReader(File file);
}
\end{javacode}
\caption{Metodi creare una risorsa \texttt{SourceResource<BufferedReader, String>}}
\label{lst:srcRefBr}
\end{elisting}
%--------1---------2---------3---------4---------5---------6---------7---------8
%--------1---------2---------3---------4---------5---------6---------7---------8
Se leggiamo un file, una riga per volta, possiamo usare il metodo statico
\texttt{bufferedReader}, sono disponibili 4 varianti, lis.~\ref{lst:srcRefBr}.
L'\,argomento \texttt{file} è obbligatorio, è possibile indicare un charset per
la lettura del file, ed è possibile indicare un'\,azione da eseguire dopo aver
letto una riga dal file (normalmente un contatore per le statistiche di
esecuzione).
%--------1---------2---------3---------4---------5---------6---------7---------8

%--------1---------2---------3---------4---------5---------6---------7---------8
\begin{elisting}[!htb]
\begin{javacode}
public interface SourceResource<U extends AutoCloseable, I> {
    static <I> SourceResource<BufferedReader, I> bufferedReader(File file, Function<String,I> dec,
                                                                Charset cs, Runnable inc);
    static <I> SourceResource<BufferedReader, I> bufferedReader(File file, Function<String,I> dec, Charset cs);
    static <I> SourceResource<BufferedReader, I> bufferedReader(File file, Function<String,I> dec, Runnable inc);
    static <I> SourceResource<BufferedReader, I> bufferedReader(File file, Function<String,I> dec);
}
\end{javacode}
\caption[Metodi creare una risorsa \texttt{SourceResource<BufferedReader,I>}
con funzione di decodifica]{Metodi creare una risorsa \texttt{SourceResource<BufferedReader,I>}
con funzione di decodifica \texttt{String}~$\mapsto$~\texttt{I}}
\label{lst:srcRefBrDec}
\end{elisting}
%--------1---------2---------3---------4---------5---------6---------7---------8
Se il file viene letto una riga per volta, e la riga viene deserializzata nella
classe dati, può essere usato un costruttore nella forma
lis.~\ref{lst:srcRefBrDec}, che sono analoghi alla forma precedente, ma hanno
la funzione di decodifica in modo che l'\,iterator o il supplier forniscano
direttamente la classe dati.
%--------1---------2---------3---------4---------5---------6---------7---------8

%--------1---------2---------3---------4---------5---------6---------7---------8
\begin{elisting}[!htb]
\begin{javacode}
public interface SourceResource<U extends AutoCloseable, I> {
    static <U extends AutoCloseable, I> SourceResource<U, I> fromIterator(Supplier<U> ctor,
                                                                          Function<U, Iterator<I>> reader);
    static <U extends AutoCloseable & Iterable<I>, I> SourceResource<U, I> fromIterator(Supplier<U> ctor);
    static <I> SourceResource<?, I> fromIterator(Iterable<I> iterable);
}
\end{javacode}
\caption{Metodi per creare una origine dati da un iterator}
\label{lst:srcRefIter}
\end{elisting}
%--------1---------2---------3---------4---------5---------6---------7---------8
Infine esistono i costruttori statici generici mediante \texttt{Iterator},
lis.~\ref{lst:srcRefIter}, e \texttt{Supplier}, lis.~\ref{lst:srcRefSupp}.
%--------1---------2---------3---------4---------5---------6---------7---------8

%--------1---------2---------3---------4---------5---------6---------7---------8
\begin{elisting}[!htb]
\begin{javacode}
public interface SourceResource<U extends AutoCloseable, I> {
    static <U extends AutoCloseable, I> SourceResource<U, I> fromSupplier(
        Supplier<U> ctor, Function<U, Supplier<I>> reader);
    static <U extends AutoCloseable & Supplier<I>, I> SourceResource<U, I> fromSupplier(Supplier<U> ctor);
}
\end{javacode}
\caption{Metodi per creare una origine dati da un supplier}
\label{lst:srcRefSupp}
\end{elisting}
%--------1---------2---------3---------4---------5---------6---------7---------8


\section{Destinazione dati --- \texttt{SinkResource}} \label{sec:snkRes}
%--------1---------2---------3---------4---------5---------6---------7---------8
Una generica destinazione dati è una classe che implementa l'\,interfaccia
\texttt{SinkResource}, vedi lis.~\ref{lst:snkRef}, dove il metodo \texttt{get}
fornisce la risorsa \texttt{U} che può essere usata con un
\texttt{try-with-resources} ed il metodo \texttt{accept} scrive\footnote{%
    uso di una \texttt{SinkResource}:
    \begin{javacode}
    try (T t = sink.get()) {    // open resource
        O o = // TODO
        sink.accept(t, o);      // write data
    }
    \end{javacode}
} il dato.
%--------1---------2---------3---------4---------5---------6---------7---------8
L'\,interfaccia oltre a indicare i metodi di istanza offre dei costruttori
statici per facilitare l'\,implementazione dell'\,interfaccia.
%--------1---------2---------3---------4---------5---------6---------7---------8

%--------1---------2---------3---------4---------5---------6---------7---------8
\begin{elisting}[!htb]
\begin{javacode}
public interface SinkResource<U extends AutoCloseable, O> {
    U get();
    void accept(U u, O o);
}
\end{javacode}
\caption{Metodi per aprire la risorsa e scrivere i dati}
\label{lst:snkRef}
\end{elisting}
%--------1---------2---------3---------4---------5---------6---------7---------8

%--------1---------2---------3---------4---------5---------6---------7---------8
Se la risorsa è un file e viene scritta una riga per volta possono essere usato
il metodo statico \texttt{bufferedWriter}, sono disponibili 4 varianti,
lis.~\ref{lst:snkRefBw}. L'\,argomento \texttt{file} è obbligatorio, è possibile
indicare il charset per la scrittura del file, e una azione da eseguire dopo
aver scritto la riga del file (normalmente un contatore per le statistiche di
esecuzione).
%--------1---------2---------3---------4---------5---------6---------7---------8

%--------1---------2---------3---------4---------5---------6---------7---------8
\begin{elisting}[!htb]
\begin{javacode}
public interface SinkResource<U extends AutoCloseable, O> {
    static SinkResource<BufferedWriter, String> bufferedWriter(File file, Charset cs, Runnable inc);
    static SinkResource<BufferedWriter, String> bufferedWriter(File file, Runnable inc);
    static SinkResource<BufferedWriter, String> bufferedWriter(File file, Charset cs);
    static SinkResource<BufferedWriter, String> bufferedWriter(File file);
}
\end{javacode}
\caption{Metodi per creare una risorsa \texttt{SinkResource<BufferedWriter, String>}}
\label{lst:snkRefBw}
\end{elisting}
%--------1---------2---------3---------4---------5---------6---------7---------8

%--------1---------2---------3---------4---------5---------6---------7---------8
Se la riga sul file rappresenta la serializzazione di un oggetto, possono essere
usati i costruttori statici nella forma lis.~\ref{lst:snkRefBwEnc}, che sono
analoghi alla forma precedente, ma hanno la funzione di codifica della classe
in stringa, in modo che il \textit{writer} possa scrivere direttamente la classe
dati.
%--------1---------2---------3---------4---------5---------6---------7---------8

%--------1---------2---------3---------4---------5---------6---------7---------8
\begin{elisting}[!htb]
\begin{javacode}
public interface SinkResource<U extends AutoCloseable, O> {
    static <O> SinkResource<BufferedWriter, O> bufferedWriter(File file, Function<O,String> enc, Charset cs,
                                                              Runnable inc);
    static <O> SinkResource<BufferedWriter, O> bufferedWriter(File file, Function<O,String> enc, Charset cs);
    static <O> SinkResource<BufferedWriter, O> bufferedWriter(File file, Function<O,String> enc, Runnable inc);
    static <O> SinkResource<BufferedWriter, O> bufferedWriter(File file, Function<O,String> enc);
}
\end{javacode}
\caption[Metodi per creare una risorsa \texttt{SinkResource<BufferedWriter,O>}
con funzione di codifica]{Metodi per creare una risorsa \texttt{SinkResource<BufferedWriter,O>}
con funzione di codifica \textit{O}~$\mapsto$~\texttt{String}}
\label{lst:snkRefBwEnc}
\end{elisting}
%--------1---------2---------3---------4---------5---------6---------7---------8

%--------1---------2---------3---------4---------5---------6---------7---------8
Infine esistono i costruttori statici generici, lis.~\ref{lst:snkRefGen}.
Il parametro \texttt{action} indica una azione che viene eseguita dopo la
scrittura del dato (normalmente un contatore per le statistiche di esecuzione).
%--------1---------2---------3---------4---------5---------6---------7---------8

%--------1---------2---------3---------4---------5---------6---------7---------8
\begin{elisting}[!htb]
\begin{javacode}
public interface SinkResource<U extends AutoCloseable, O> {
    static <U extends AutoCloseable, O> SinkResource<U, O> of(Supplier<U> ctor, BiConsumer<U, O> writer);
    static <U extends AutoCloseable, O> SinkResource<U, O> of(Supplier<U> ctor, BiConsumer<U, O> writer,
                                                              Runnable action);
    static <O> SinkResource<?, O> of(Consumer<O> writer);
    static <O> SinkResource<?, O> of(Consumer<O> writer, Runnable action);
}
\end{javacode}
\caption{Metodi per creare una risorsa \texttt{SinkResource} generica}
\label{lst:snkRefGen}
\end{elisting}
%--------1---------2---------3---------4---------5---------6---------7---------8

\subsection{Destinazione dati --- \texttt{SinkBufResource}} \label{sec:snkBufRes}
%--------1---------2---------3---------4---------5---------6---------7---------8
A volte una destinazione dati non \textit{scrive} immediatamente i dati che
riceve, ma gli accumula in un buffer, e la scrittura avviene al raggiungimento
di un limite del buffer o prima della chiusura della risorsa.
Questa implementazione in generale è più performante rispetto alla scrittura
immediata.

%--------1---------2---------3---------4---------5---------6---------7---------8
Quando si hanno più destinazioni dati, in generale, le soglie dei rispettivi
buffer non vengono raggiunte in modo sincrono, e i dati vengono scritti,
temporaneamente, in modo inconsistente.
Questo non è un problema se le risorse sono files, ma se le risorse sono update
(insert) batch su tabelle diverse, collegate, di un database, potrebbe essere
un problema.

%--------1---------2---------3---------4---------5---------6---------7---------8
L'\,interfaccia \texttt{SinkBufResource} è banalmente una estensione della
interfaccia \texttt{SinkResource}, che rispetto a questa, ha in più il metodo
\texttt{flush()} che permette di forzare la scritura del buffer.
Questo permette di impostare i legami delle risorse e scrivere i dati sulle
risorse collegate in modo consistente.

%--------1---------2---------3---------4---------5---------6---------7---------8
\begin{elisting}[!htb]
\begin{javacode}
public interface SinkBufResource<U extends AutoCloseable, O> extends SinkResource<U, O> {
	void flush();
	static <U extends AutoCloseable, O>
    SinkBufResource<U, O> of(Supplier<U> ctor, BiConsumer<U, O> writer, Consumer<U> flusher);
    static <U extends AutoCloseable, O>
    SinkBufResource<U, O> of(Supplier<U> ctor, BiConsumer<U, O> writer, Consumer<U> flusher, Runnable action);
	...
}
\end{javacode}
\caption{Metodi per creare una risorsa \texttt{SinkBufResource} generica}
\label{lst:snkBufResGen}
\end{elisting}
%--------1---------2---------3---------4---------5---------6---------7---------8

%--------1---------2---------3---------4---------5---------6---------7---------8
\begin{elisting}[!htb]
\begin{javacode}
public interface SinkBufResource<U extends AutoCloseable, O> extends SinkResource<U, O> {
	...
    static SinkBufResource<BufferedWriter, String> bufferedWriter(File file, Charset cs, Runnable inc);
    static SinkBufResource<BufferedWriter, String> bufferedWriter(File file, Runnable inc);
    static SinkBufResource<BufferedWriter, String> bufferedWriter(File file, Charset cs);
    static SinkBufResource<BufferedWriter, String> bufferedWriter(File file);
	...
}
\end{javacode}
\caption{Metodi per creare una risorsa \texttt{SinkBufResource<BufferedWriter, String>}}
\label{lst:snkBufResStr}
\end{elisting}
%--------1---------2---------3---------4---------5---------6---------7---------8

%--------1---------2---------3---------4---------5---------6---------7---------8
\begin{elisting}[!htb]
\begin{javacode}
public interface SinkBufResource<U extends AutoCloseable, O> extends SinkResource<U, O> {
	...
    static <O> SinkBufResource<BufferedWriter, O> bufferedWriter(
    		File file, Function<O, String> enc, Charset cs, Runnable inc);
    static <O> SinkBufResource<BufferedWriter, O> bufferedWriter(
    		File file, Function<O, String> enc, Runnable inc);
    static <O> SinkBufResource<BufferedWriter, O> bufferedWriter(File file, Function<O, String> enc, Charset cs);
    static <O> SinkBufResource<BufferedWriter, O> bufferedWriter(File file, Function<O, String> enc);
}
\end{javacode}
\caption{Metodi per creare una risorsa \texttt{SinkBufResource<BufferedWriter, O>} con funzione di codifica}
\label{lst:snkBufResFcn}
\end{elisting}
%--------1---------2---------3---------4---------5---------6---------7---------8

